% ==============================================================================
% Capítulo 4: Metodologia
% ==============================================================================
\chapter{Metodologia}
\label{cap:metodologia}

\begin{instrucao}
\orientacao{O que vai na Metodologia}
Este capítulo descreve o que foi feito e como, com detalhes suficientes para
que outra pessoa consiga reproduzir o trabalho. Em um trabalho científico de
Computação, ele costuma ter duas partes:
\begin{itemize}
  \item a \textbf{abordagem proposta}: modelo, algoritmo, arquitetura ou
        técnica desenvolvida, com as decisões de projeto justificadas;
  \item o \textbf{planejamento dos experimentos}: dados ou instâncias,
        métricas, ambiente computacional, parâmetros e métodos de comparação.
\end{itemize}
Escreva no passado e na voz passiva ou impessoal (``a heurística foi
implementada'', ``foram usadas 270 instâncias''). Os resultados ficam para o
capítulo seguinte. Aqui apenas se descreve como eles foram obtidos.

Se o título ``Metodologia'' não combinar com o trabalho, ele pode ser trocado
pelo nome da abordagem, como ``Heurística proposta'' ou ``Método proposto''.
\end{instrucao}

Este capítulo apresenta a heurística proposta para o \acs{mkp} e o
planejamento dos experimentos usados para avaliá-la. A
\cref{sec:metodologia:heuristica} descreve os componentes da heurística e a
\cref{sec:metodologia:experimentos} apresenta as instâncias, as métricas e o
ambiente computacional.

% ------------------------------------------------------------------------------
\section{Heurística proposta}
\label{sec:metodologia:heuristica}

A heurística segue o esquema da \acs{ils} apresentado na
\cref{sec:fundamentacao:ils} e está descrita no \cref{alg:ils}. As subseções a
seguir detalham cada componente.

\begin{algoritmo}[htbp]
  \caption{Heurística de \acs{ils} proposta para o \acs{mkp}}
  \label{alg:ils}
  \begin{algorithmic}[1]
    \Require instância $(p, w, c)$, tempo limite $T$, parâmetros $k_{\min}$ e $k_{\max}$
    \Ensure melhor solução encontrada $s^{\mathrm{best}}$
    \State $s \gets \Call{Gulosa}{p, w, c}$
    \State $s \gets \Call{BuscaLocal}{s}$
    \State $s^{\mathrm{best}} \gets s$
    \State $k \gets k_{\min}$
    \While{tempo decorrido $< T$}
      \State $s' \gets \Call{Perturba}{s, k}$ \Comment{remove $k$ itens ao acaso}
      \State $s' \gets \Call{BuscaLocal}{s'}$
      \If{$f(s') > f(s)$}
        \State $s \gets s'$
        \State $k \gets k_{\min}$
        \If{$f(s) > f(s^{\mathrm{best}})$}
          \State $s^{\mathrm{best}} \gets s$
        \EndIf
      \Else
        \State $k \gets \min(k + 1, k_{\max})$ \Comment{intensifica a perturbação}
      \EndIf
    \EndWhile
    \State \Return $s^{\mathrm{best}}$
  \end{algorithmic}
  \fonteautor
\end{algoritmo}

\begin{instrucao}
\orientacao{Algoritmos em pseudocódigo}
Algoritmos são apresentados em pseudocódigo, e não como código-fonte de uma
linguagem específica. O modelo usa o pacote \verb|algpseudocode| dentro do
ambiente \verb|algoritmo|, que tem lista própria e legenda acima, como as
demais ilustrações.
\begin{exemplocodigo}
\begin{algoritmo}[htbp]
  \caption{Busca linear}
  \label{alg:busca}
  \begin{algorithmic}[1]
    \Require vetor $A[1..n]$, valor $x$
    \Ensure posição de $x$ em $A$ ou $-1$
    \For{$i \gets 1$ \textbf{até} $n$}
      \If{$A[i] = x$}
        \State \Return $i$
      \EndIf
    \EndFor
    \State \Return $-1$
  \end{algorithmic}
  \fonteautor
\end{algoritmo}
\end{exemplocodigo}
O número entre colchetes em \verb|algorithmic[1]| numera todas as linhas, o
que permite citar uma linha específica no texto (``a linha 6 do
\cref{alg:ils}''). Explique o algoritmo no texto, linha a linha ou por blocos,
como nas subseções abaixo.
\end{instrucao}

\subsection{Solução inicial}
\label{sec:metodologia:inicial}

A solução inicial (linha 1 do \cref{alg:ils}) é construída por um procedimento
guloso que ordena os itens pela razão de eficiência
\begin{equation}
  e_j = \frac{p_j}{\sum_{i=1}^{m} w_{ij} / c_i},
  \label{eq:eficiencia}
\end{equation}
em que o consumo de cada recurso é ponderado pela respectiva capacidade. Os
itens são inseridos em ordem decrescente de $e_j$, desde que a inserção não
viole nenhuma restrição.

\subsection{Busca local}
\label{sec:metodologia:busca}

A busca local explora duas vizinhanças. A primeira tenta inserir qualquer item
fora da solução que caiba na capacidade residual. A segunda troca um item da
solução por um item de fora, desde que a troca aumente o valor da solução e
mantenha a viabilidade. A busca adota a estratégia de primeira melhora e
termina quando nenhuma das duas vizinhanças produz uma solução melhor.

\subsection{Perturbação adaptativa e critério de aceitação}
\label{sec:metodologia:perturbacao}

A perturbação remove $k$ itens escolhidos ao acaso da solução corrente. O valor
de $k$ começa em $k_{\min}$ e aumenta em uma unidade a cada iteração sem
melhora, até o limite $k_{\max}$ (linha 15). Quando uma solução melhor é
encontrada, $k$ volta a $k_{\min}$ (linha 10). Desse modo, a busca começa com
perturbações fracas, que exploram a vizinhança do ótimo local, e as intensifica
apenas quando fica estagnada. O critério de aceitação é o de melhora, ou seja,
a busca continua a partir de $s'$ somente se $s'$ for melhor que $s$.

% ------------------------------------------------------------------------------
\section{Planejamento dos experimentos}
\label{sec:metodologia:experimentos}

\subsection{Instâncias e métodos de comparação}

Foram usadas as 270 instâncias de \textcite{chu1998genetic}, organizadas em
27 grupos de dez instâncias que combinam $n \in \{100, 250, 500\}$ itens,
$m \in \{5, 10, 30\}$ restrições e três níveis de folga das restrições. A
heurística foi comparada ao modelo das
\crefrange{eq:mkp:objetivo}{eq:mkp:dominio}, resolvido pelo HiGHS
\cite{huangfu2018parallelizing}, e aos resultados do \acs{ag} de
\textcite{chu1998genetic}.

\subsection{Métricas}

A qualidade de uma solução de valor $z$ foi medida pelo desvio percentual
(\emph{gap}) em relação ao valor $z^*$ da melhor solução conhecida para a
instância, dado pela \cref{eq:gap}.
\begin{equation}
  \mathit{gap} = 100 \times \frac{z^* - z}{z^*}.
  \label{eq:gap}
\end{equation}
Como a heurística é aleatória, cada instância foi resolvida dez vezes com
sementes diferentes e foram registrados o \emph{gap} médio e o melhor
\emph{gap}. A diferença entre os métodos foi avaliada pelo teste de postos
sinalizados de Wilcoxon, com nível de significância de 5\%.

\subsection{Ambiente computacional e parâmetros}

A heurística foi implementada em C++17 e compilada com o g++ 13 e a opção
\texttt{-O3}. Os experimentos foram executados em um computador com processador
Intel Core i7 de 12ª geração, 16 GB de memória e sistema Ubuntu 24.04. Todos
os métodos receberam tempo limite de 60 segundos por instância. Os parâmetros
$k_{\min} = 2$ e $k_{\max} = 10$ foram definidos em testes preliminares com 27
instâncias, uma de cada grupo, que não foram usadas na comparação final.

\begin{instrucao}
\orientacao{Reprodutibilidade}
Informe tudo o que influencia os resultados, como a versão das ferramentas,
o hardware, as sementes aleatórias, os parâmetros e a forma de calibrá-los. Uma
boa prática é publicar o código e os dados em um repositório (GitHub, GitLab
ou Zenodo) e indicar o endereço no texto ou em nota de rodapé.
\end{instrucao}

% ------------------------------------------------------------------------------
\section{Uso de ferramentas de inteligência artificial}
\label{sec:metodologia:ia}

Durante o trabalho, o ChatGPT \cite{openai2026chatgpt}, com o modelo GPT-5 e
as configurações padrão da interface web, foi usado para revisar a redação de
trechos do texto e para sugerir casos de teste para a rotina que verifica a
viabilidade das soluções. Nenhum dado pessoal ou confidencial foi fornecido à
ferramenta. Todo conteúdo sugerido foi conferido pelo autor antes de ser usado.
A declaração completa e o detalhamento das interações estão no \cref{ap:ia}.

\begin{instrucao}
\orientacao{Declaração de uso de inteligência artificial}
O Art. 19 da Portaria Normativa nº 310/2026 da Ufes \cite{ufes2026portaria}
torna obrigatória a declaração do uso de \ac{ia} em qualquer produção
acadêmica, exceto os usos meramente acessórios definidos pela comissão de
governança de IA da Ufes (CPGIA). Para trabalhos acadêmicos, a portaria pede
três coisas. O modelo atende a cada uma delas em um lugar:
\begin{enumerate}
  \item uma declaração com ferramenta, versão, finalidade e forma de
        contribuição, em seção ou apêndice específico (§ 1º): o
        \cref{ap:ia}, que fica como primeiro apêndice para ser encontrado com
        facilidade;
  \item o detalhamento das interações, das motivações e dos métodos de uso
        (§ 2º): também no \cref{ap:ia};
  \item a citação da ferramenta conforme as normas de citação e a indicação,
        na descrição dos métodos, da versão, dos parâmetros e das bases de
        dados usadas (§ 3º): esta seção, com a ferramenta nas Referências.
\end{enumerate}
A portaria também proíbe atribuir autoria ou coautoria a uma ferramenta de IA
(Art. 17) e inserir nela dados pessoais ou confidenciais. Você continua
responsável pela correção e pela originalidade de tudo o que entregar
(Art. 18). Se não usou nenhuma ferramenta de IA, apague esta seção e o
\cref{ap:ia}. Quando o uso decorrer de tecnologia assistiva, a declaração pode
ser apresentada ao setor de apoio à inclusão, e não no corpo do trabalho
(§ 4º). O modelo oficial de declaração ainda será publicado pela CPGIA. Quando
isso acontecer, ajuste o apêndice a ele.
\end{instrucao}
